% BAB 1 - PENDAHULUAN

\vspace{1cm}
\subsection{Deskripsi Aplikasi}
\indent
Perkembangan teknologi informasi mendorong usaha kecil, termasuk laundry, untuk beralih dari pencatatan konvensional ke sistem digital. Salah satu solusi yang dikembangkan adalah Sistem Point of Sale (POS) Laundry berbasis web, yang berfungsi untuk mencatat transaksi, mengelola data pelanggan, memantau status cucian, serta menyusun laporan operasional secara terintegrasi. Dengan berbasis web, sistem ini dapat diakses secara fleksibel dari berbagai perangkat, sehingga meningkatkan efisiensi dan memudahkan pemilik usaha dalam melakukan monitoring \cite{Achyani2024}.

Penelitian sebelumnya menegaskan bahwa aplikasi laundry berbasis web mampu menyederhanakan proses pencarian dan pengelolaan layanan laundry, sekaligus meningkatkan transparansi data dan akurasi pencatatan \cite{Singh2022}. Selain itu, pendekatan berbasis pengguna dalam pengembangan sistem laundry terbukti dapat memperbaiki alur operasional dan meningkatkan kepuasan pelanggan, terutama melalui antarmuka yang mudah digunakan dan proses transaksi yang lebih cepat \cite{Kalo2023}.

Studi lokal juga menunjukkan bahwa sistem informasi laundry berbasis web dapat membantu pemilik usaha dalam mengurangi kesalahan pencatatan, mempercepat proses transaksi, serta menghasilkan laporan yang lebih rapi dan terstruktur \cite{SigitSugiharto2023}. Penelitian lain menekankan bahwa penerapan sistem digital pada usaha kecil menengah (UMKM) berpengaruh positif terhadap efisiensi operasional dan kualitas layanan, sehingga mendukung keberlanjutan bisnis \cite{Maulana2025}.

Berdasarkan kondisi tersebut, usaha laundry yang masih bergantung pada pencatatan berbasis kertas berisiko menghadapi kendala seperti kesalahan hitung, keterlambatan informasi, dan pengelolaan data yang kurang optimal. Oleh karena itu, pengembangan Sistem POS Laundry Berbasis Website diharapkan menjadi solusi yang mampu meningkatkan efisiensi operasional, kualitas layanan, serta kepuasan pelanggan secara berkelanjutan.

\subsection{Identifikasi Masalah}
\noindent
Berdasarkan latar belakang yang telah dijelaskan, maka pertanyaan penelitian (\textit{research questions}) yang diidentifikasi dalam penelitian ini adalah sebagai berikut:

\begin{enumerate}
    \item Bagaimana merancang sistem POS Laundry berbasis web yang mampu mentransformasi pencatatan konvensional menjadi sistem informasi yang terintegrasi?
    \item Bagaimana mengimplementasikan mekanisme \textit{callback} notifikasi pembayaran QRIS untuk mengautomasi verifikasi transaksi secara \textit{real-time}?
    \item Bagaimana menjamin keandalan fitur pembayaran digital pada sistem menggunakan pengujian unit dengan target 100\% \textit{statement coverage}?
    \item Bagaimana penerapan metode Scrum dapat mendukung pengembangan sistem yang adaptif terhadap kebutuhan operasional UMKM Laundry?
\end{enumerate}

\noindent
Pertanyaan penelitian ini dirumuskan untuk menjawab kebutuhan akan sistem yang tidak hanya mendigitalisasi data, tetapi juga membawa pembaruan pada aspek automasi dan validasi data transaksi \cite{Achyani2024,Mutia2025}.

\subsection{Tujuan}
\noindent
Tujuan dari penelitian dan pengembangan sistem POS Laundry berbasis web ini adalah:

\begin{enumerate}
    \item Merancang dan membangun sistem POS Laundry berbasis web yang terintegrasi untuk meningkatkan efisiensi operasional.
    \item Membangun fitur automasi verifikasi pembayaran QRIS yang mampu mendeteksi dana masuk secara otomatis guna meminimalkan celah manipulasi data.
    \item Mencapai tingkat \textit{statement coverage} sebesar 100\% pada modul krusial untuk memastikan sistem bebas dari kesalahan logika.
    \item Mengoptimalkan pengelolaan data pelanggan dan laporan keuangan melalui basis data yang terstruktur.
\end{enumerate}

\noindent
Tujuan ini didukung oleh penelitian sebelumnya yang menunjukkan bahwa sistem berbasis web dapat meningkatkan efisiensi operasional, memperbaiki kualitas layanan, dan mendukung keberlanjutan usaha laundry \cite{Achyani2024,Kalo2023,SigitSugiharto2023,Maulana2025}.

\subsection{Lingkup Dokumentasi}
\noindent
Lingkup dokumentasi pada proyek ini dibatasi pada aspek-aspek utama yang mendukung pengembangan Sistem POS Laundry berbasis web, meliputi:

\begin{enumerate}
    \item Analisis kebutuhan sistem yang mencakup identifikasi modul inti seperti transaksi, pelanggan, dan laporan.
    \item Perancangan sistem berbasis web yang meliputi desain antarmuka sederhana, responsif, dan mudah digunakan.
    \item Implementasi modul inti berupa transaksi laundry, manajemen pelanggan, serta laporan operasional dan keuangan.
    \item Evaluasi awal sistem melalui uji coba terbatas untuk menilai efisiensi dan akurasi pencatatan.
\end{enumerate}

\noindent
Batasan dokumentasi tidak mencakup pengembangan aplikasi mobile, integrasi notifikasi otomatis, maupun fitur tambahan di luar modul inti. Fokus diarahkan pada sistem berbasis web yang mendukung pencatatan transaksi, pemantauan status cucian, dan penyusunan laporan secara terstruktur \cite{Mutia2025,Lizal2025,Kalo2023}.
